\section{Small fourfold detectors after Schoen and van Geemen}
\label{sec:small}
This section computes degrees and integral coefficients for the geometric cycles of Schoen and van Geemen. The construction of the cycles is their input; the estimates keep track of the polarization, the isogeny and the saturated Weil lattice.

\begin{theorem}\label{thm:small-detectors}
Let $(A,\theta)$ be a principally polarized Weil fourfold with full integral maximal-order action. For $K=\Q(i)$ the detector degree is at most $48$ on the even Hermitian lattice $U\oplus U$, and at most $192$ on the odd lattice $I_{2,2}$. For $K=\Q(\sqrt{-3})$ it is at most $36$. In the last case, if $\mathcal W=W_K(A)\cap H^4(A,\Z)$, then $3\mathcal W$ consists of integral algebraic cycle classes. For $0\ne t\in W_K(A)$, $\delta\in\Z_{>0}$ with $\delta t$ integral, and $T=\int_A t^2$, one has
\begin{equation}\label{eq:eisenstein-cost}
 c_t\leq\max\left\{3\delta,
 36\left\lfloor\frac{\delta\sqrt{2T}}3\right\rfloor\right\}
 \leq17\delta\lceil\sqrt T\rceil.
\end{equation}
Here $U=\left(\begin{smallmatrix}0&1\\1&0\end{smallmatrix}\right)$ and $I_{r,s}$ is the diagonal Hermitian form with $r$ entries $1$ and $s$ entries $-1$.
\end{theorem}

\subsection{The Gaussian bound in two lattice types}

Van Geemen \cite[Proposition~3.6 and Theorem~3.7, pp.~623--625]{vG96} constructs a dominant rational map from a general member of the even principal Gaussian family to a smooth quadric fourfold in $\mathbf P^5$. The six defining sections belong to $|2\theta|$ and have finite base locus; the inverse image of a suitable ruling plane has nonzero Weil projection. This yields the degree bound with a short residual-intersection argument.

Choose such a plane $\Pi$ generally in its ruling, and choose a linear $\mathbf P^3$ through it. Its intersection with the quadric is the union of $\Pi$ and a residual ruling plane. The two pulled-back hyperplanes have no common divisor for general choices: away from the finite base locus this follows by avoiding the finitely many divisorial components contracted by the rational map. Their proper intersection is an effective surface cycle of class $4\theta^2$ and degree
\[
 \int_A4\theta^4=96.
\]
It contains a detecting component $Y$ from the inverse image of $\Pi$. If $\deg Y\leq48$, this is the desired detector. Otherwise subtract one copy of $Y$ from the complete-intersection cycle. The residual effective cycle has degree less than $48$ and has nonzero Weil projection, since the full class $4\theta^2$ has zero Weil projection. One of its components detects. Proposition~\ref{prop:detector-loci} extends the bound from the general member to every member of the even component.

For the odd lattice put
\[
 P=\begin{pmatrix}1&1\\1&-1\end{pmatrix},\qquad F=P\oplus P.
\]
Direct multiplication gives $P^*UP=2I_{1,1}$ and $\Nm(\det F)=16$. Transporting the complex structure by $F$ gives a $K$-linear isogeny
\[
 f:A_{\rm odd}\longrightarrow A_{\rm even},\qquad
 \deg f=16,\quad f^*\theta_{\rm even}=2\theta_{\rm odd}.
\]
For an effective surface cycle $Z$ on the target,
\[
 \deg_{\theta_{\rm odd}} f^*Z
 =\frac{16}{2^2}\deg_{\theta_{\rm even}}Z.
\]
Pullback preserves nonzero Weil projection, because it is a $K$-linear cohomology isomorphism. Pulling back an even detector and taking a detecting component proves the bound $192$.

Combining these bounds with Proposition~\ref{prop:gaussian-extraction} gives, respectively,
\[
 c_t\leq141\,557\,760\,\delta\lceil\sqrt T\rceil,
 \qquad
 c_t\leq2\,264\,924\,160\,\delta\lceil\sqrt T\rceil.
\]

\subsection{The Schoen surface and its principal quotient}

The published geometric input is the deformation of
$\Delta_J\cup(C\times C)\subset J\times J$, for a genus-two Jacobian $J$, to an embedded Schoen surface in a four-dimensional Eisenstein Weil family. We use the form of Schoen's construction \cite{Schoen07} stated by van Geemen \cite[Introduction, p.~1, and \S6.2, pp.~14--15]{vG22}, including the ambient family and its polarization. The embedded flat degeneration is also described in \cite[\S3, pp.~5--6]{Beauville}, and the four-dimensional smooth surface moduli in \cite[Theorem~1]{RRS}.

For clarity, the dimension of the surface moduli really gives dominance onto the Weil component. A smooth Schoen surface $S\subset A$ is Lagrangian for a nondegenerate holomorphic two-form and has $q(S)=4$ \cite[Introduction]{RRS}. Thus $N_{S/A}\simeq\Omega_S^1$ and $h^0(N_{S/A})=4$. Translation supplies four independent infinitesimal embedded deformations: the infinitesimal stabilizer is zero because a smooth surface of general type has no vector fields. Therefore a fixed abelian variety admits no positive-dimensional family of these embedded surfaces modulo translation. The map from the four-dimensional surface moduli to Albanese moduli has zero-dimensional fibres. Its image in the four-dimensional Weil component is dominant. This supplies an analytic open set of the component carrying the required effective surface class.

We compute its degree and Weil part at the degeneration. Let $\vartheta$ be the principal class on $J$, and on $J^2$ set
\[
 x=p_1^*\vartheta,\qquad y=p_2^*\vartheta,\qquad
 z=m^*\vartheta-x-y.
\]
The polarization $a x+b z+c y$ corresponds to the matrix
$\left(\begin{smallmatrix}a&b\\b&c\end{smallmatrix}\right)$ on the two factors. The determinant formula for its Euler characteristic, or expansion in a symplectic basis, gives
\begin{equation}\label{eq:intersection-polynomial}
 \int_{J^2}(a x+b z+c y)^4=24(ac-b^2)^2.
\end{equation}
In particular
\[
 \int x^2y^2=4,\qquad \int xyz^2=-4,\qquad
 \int z^4=24,
\]
and the other degree-four monomial intersections vanish. The central surface has class
\[
 s=[\Delta_J]+[C\times C]
   =\tfrac12(x+y-z)^2+xy.
\]
Take $\omega=(-1+\sqrt{-3})/2$ and the endomorphism and polarization
\[
 R=\begin{pmatrix}0&-1\\1&-1\end{pmatrix},\qquad
 H=2x-z+2y.
\]
The matrix of $H$ is $M=\left(\begin{smallmatrix}2&-1\\-1&2\end{smallmatrix}\right)$ and $R^tMR=M$. Its polarization type is $(1,1,3,3)$ and $H^4=216$. Equation~\eqref{eq:intersection-polynomial} gives
\[
 \int H^2s=36,\qquad \int s^2=8.
\]
For example, the two surface components contribute respectively $8$ and $28$ to the first number. The diagonal has self-intersection zero, the product has self-intersection $4$, and their mutual intersection is $2$, giving the second number.

Put $\ell_+=\omega^2x+\omega z+y$ and
$\ell_-=\omega x+\omega^2z+y$. Their squares span the two extreme $K$-eigenlines. Direct expansion gives
\[
 w_0=s-\tfrac16H^2
   =\tfrac16\bigl(\omega\ell_+^2+\omega^2\ell_-^2\bigr)\in W,
 \qquad \int w_0^2=2.
\]
The same flat class and norm hold on the deformed surface.

The polarization kernel is $\{(u,-u):u\in J[3]\}$, and $R$ acts trivially on it. A maximal isotropic subgroup of order $9$ therefore gives an Eisenstein-linear quotient
\[
 q:A_H\longrightarrow A,\qquad \deg q=9,\qquad q^*\theta=H,
\]
with principal polarization on $A$. The quotient can be chosen over a finite cover of the period family. To identify its principal lattice, specialize $J$ to $E_\omega^2$. The explicit map
\[
 F=\begin{pmatrix}1&\omega\\1&\omega^2\end{pmatrix}
       \otimes I_{E_\omega^2}
\]
has degree $9$, satisfies $F^*\theta_{E_\omega^4}=H$, and intertwines $R$ with $\operatorname{diag}(\omega,\omega^2)$. Its target is the balanced product $E_\omega^4$ with product polarization. Transporting this finite-index lattice along the period domain identifies the whole principal component.

This component covers every principal Eisenstein lattice of signature $(2,2)$. Indeed the maximal order is a principal ideal domain, and the associated integral Hermitian lattice is unimodular with determinant $1$. It is rationally split by the classification in \cite[Proposition~4.1 and Corollary~4.2]{Deligne82}. Choose a primitive integral isotropic vector $v$. Unimodularity gives $w$ with $h(v,w)=1$. If $h(w,w)=b\in\Z$, replacing $w$ by $w+b\omega v$ makes it isotropic because $\operatorname{Tr}(\omega)=-1$. Thus $v,w$ span a unimodular hyperbolic plane, which splits off integrally. Repeat on its signature-$(1,1)$ complement. The lattice is $U\oplus U$. Finally, in $I_{1,1}$ the vectors $(1,1)$ and $(1+\omega,\omega)$ have Gram matrix $U$ and determinant $-1$, so $U\oplus U\simeq I_{2,2}$.

Push the entire effective surface cycle forward, including its multiplicities, and write $Z=q_*S$. Then
\begin{equation}\label{eq:degree36-class}
 \deg_\theta Z=36,\qquad
 [Z]=\tfrac32\theta^2+w,\qquad
 w\in\mathcal W,\qquad \int_Aw^2=18.
\end{equation}
The degree follows by projection formula. For the norm use $q^*q_*s=9s$ in cohomology, so $(q_*s)^2=9s^2=72$, and subtract $(3\theta^2/2)^2=54$. The integrality of $w$ follows from that of $\theta^2/2$, which is immediate by expanding the square of an integral alternating two-form. Proper cycle parameter spaces extend an effective cycle with these same flat relations from the nonempty open set to every fibre: the components on which these relations hold are open and closed, and at least one dominates the base. Their proper image is all the base. This proves the degree-$36$ detector bound.

\subsection{The saturated Eisenstein lattice}

It remains to turn the effective cycle into the integral coefficient bound. Compute at the balanced product $E_\omega^4$. Choose integral one-form bases $x_j,y_j$, adapted to the maximal-order action, and write
\[
 \prod_{j=1}^4(x_j+\omega y_j)=u+\omega v.
\]
The all-$x$ coefficient of $(u,v)$ is $(1,0)$, and any coefficient with exactly one $y$ is $(0,1)$. Consequently a rational combination of $u,v$ is integral only when both coefficients are integers. They are a basis of the saturated Weil lattice. The identities
\[
 (u+\omega v)^2=0,\qquad
 \int(u+\omega v)(u+\omega^2v)=9
\]
give
\begin{equation}\label{eq:eisenstein-gram}
 \operatorname{Gram}(\mathcal W)=
 \begin{pmatrix}6&3\\3&6\end{pmatrix}.
\end{equation}
These are integral local systems, so the result holds throughout the family. If $\rho$ is the automorphism induced by $\omega$, then
\[
 \rho^*u=-v,\qquad \rho^*v=u-v.
\]
The norm-$18$ vectors in this lattice have coordinates
$\pm(1,1),\pm(2,-1),\pm(-1,2)$: this follows by solving
$m^2+mn+n^2=3$. The three unit translates of any such vector have differences giving, up to signs and a lattice isometry, the three directions $3u,3v,3(u-v)$. Applied to the Weil part in \eqref{eq:degree36-class}, the differences of the three degree-$36$ cycles
$Z,\rho^*Z,\rho^{2*}Z$ therefore generate exactly $3\mathcal W$.

Write $\delta t=mu+nv$. It is a sum of
\[
 N=\max\{|m|,|n|,|m+n|\}
\]
signed steps in the directions $u,v,u-v$: for equal signs use $u$ and $v$, and for opposite signs first use $u-v$ and then the remaining direction. Put $Q=m^2+mn+n^2=\delta^2T/6$. The three identities
\[
 4Q-3m^2=(m+2n)^2,\quad
 4Q-3n^2=(2m+n)^2,\quad
 4Q-3(m+n)^2=(m-n)^2
\]
show that $N\leq\delta\sqrt{2T}/3$. Each step, multiplied by $3$, is represented by a difference of two degree-$36$ effective cycles. Hence $3\delta t$ has effective numerator degrees at most $36N$. This proves the first inequality in \eqref{eq:eisenstein-cost}; the second follows from $12\sqrt2<17$ and $\lceil\sqrt T\rceil\geq1$.

\begin{example}
For the integral class $t=u$ in \eqref{eq:eisenstein-gram}, one has $T=6$ and may take $\delta=1$. A single unit-orbit difference represents $3u$, with both effective degrees $36$. The floor in \eqref{eq:eisenstein-cost} is then $\lfloor\sqrt{12}/3\rfloor=1$, exactly reproducing this presentation.
\end{example}
